
% Sprach- und Zeichensatzeinstellungen
\usepackage[english]{babel}  
\usepackage{pdfpages}
% Mehrsprachigkeit (Deutsch und Englisch)
\usepackage[T1]{fontenc}            % Zeichensatz
\usepackage[utf8]{inputenc}         % UTF-8 Eingabe

% Seitenlayout
\usepackage[a4paper, left=2cm, right=2cm, top=2.5cm, bottom=2.5cm]{geometry}

% Kopf- und Fußzeile
\usepackage{fancyhdr} 
\pagestyle{fancy}
\fancyhf{} % Kopf- und Fußzeilen leeren
\fancyfoot[C]{\thepage}

% Linien für die Kopfzeile deaktivieren
%\renewcommand{\headrulewidth}{1 pt} 
%%%%%%%%%% Wegen Format bashing we need to stick to one-side
%\setlength{\headheight}{35pt}
%\fancyhead[L]{%
%    \ifodd\value{page}
%        {\large \textit{Düsseldorf Student Economic Review Vol. 2}}\\
%        \vspace*{-0.3cm}
%        
%    \else
%        {}
%    \fi
%}

%\fancyhead[R]{%
%    \ifodd\value{page}
%        {}
%    \else
%        {\large \textit{\articletitle}}\\
%        \vspace*{-0.3cm}
%    \fi
%}



% Grafiken und Tabellen
\usepackage{graphicx}                 % Einfügen von Bildern


\graphicspath{{../figures/Base_Case/}{../figures/Dynamic_MisSpec/}{../figures/NonLinear_DGP/}} %Ref figures outside Latex Directory
\usepackage{caption}                  % Bildunterschriften
\usepackage{tabularx}                   % Einfügen von Tabellen
\usepackage{subcaption}               % Mehrere Bilder nebeneinander
\usepackage{booktabs}                 % Tabellenlinien
\usepackage{threeparttable}           % Tabellenfußnoten
\usepackage{float}                    % Flexible Positionierung von Objekten

% Farben und Hyperlinks
\usepackage[hidelinks]{hyperref}      % Hyperlinks ohne sichtbare Rahmen
\usepackage{xcolor}                   % Farben für Text und Elemente
\usepackage{colortbl}               % Für Farben in Tabellen

% Mathematik
\usepackage{amsmath, amsfonts, mathrsfs} % Erweiterte Mathematikpakete
\usepackage{bbm} %math numbers
\usepackage{dsfont}
\usepackage{bbold} % coole 1 

% Algorithmen und Code
\usepackage{algorithm}                % Algorithmenumgebung
\usepackage{algpseudocode}            % Pseudocode für Algorithmen
\usepackage{listings}                 % Quellcode-Darstellung

% Nummerierte und unnummerierte Listen
\usepackage{enumitem}                 % Anpassbare Listen

% Struktur und Titel/Verzeichnisse
\usepackage{titlesec}                 % Anpassung von Überschriften
\titleformat{\section}                % Redefinition der Sections
  {\normalfont\Large\bfseries}{\thesection.}{1em}{}
  \usepackage{titlesec}
\setcounter{tocdepth}{4} % Tiefe des Inhaltsverzeichnisses
\usepackage[backend=biber, style=apa, maxcitenames=1, mincitenames=1, uniquename=false, uniquelist = false, labeldateparts=true]{biblatex}  % Literaturverzeichnis
% Bei mehr als zwei Autoren: Zeige nur den ersten Autor mit "et al."
%\AtEveryCitekey{\ifthenelse{\value{listtotal} > 2}{\defcounter{maxnames}{1}}{}}

\usepackage{url} % damit Urls im Literaturverzeichnis nicht die Seitenränder überschreiten  
\setcounter{biburlnumpenalty}{100}
\setcounter{biburlucpenalty}{100}
\setcounter{biburllcpenalty}{100}

\DeclareDelimFormat{finalnamedelim}{\addspace\&\space}
\DeclareDelimFormat{multinamedelim}{\addcomma\space}



% Customize page prefixes for each language
\DefineBibliographyStrings{english}{
  page = {p.}, % Prefix for English pages
  pages = {pp.} % Prefix for multiple pages in English
}

% einzelne Literaturverzeichnisse
\addbibresource{bibtex.bib}

% Zusätzliche Pakete für Spezialfälle
\usepackage{authblk}                  % Autoreninformationen
\usepackage{appendix}                 % Anhänge
\usepackage{csquotes}                 % Zitate
\usepackage{diagbox}                  % Diagonale Zellen in Tabellen
\usepackage{tcolorbox}                % Farbboxen für Hervorhebungen
\usepackage{tikz}                     % Grafiken
\usepackage{wrapfig}
\usetikzlibrary{arrows.meta, positioning} % TikZ-Erweiterungen


% Tabellen-Unterschriften oben 
\captionsetup[table]{position=top}   % Tabellenunterschrift oben

% Abstände und Zeilenabstand
%\usepackage{parskip}
%\setlength{\parskip}{1em}             % Abstand zwischen Absätzen
\usepackage[onehalfspacing]{setspace} % 1,5 Zeilenabstand


% Sonderzeichen
\usepackage{eurosym}

\usepackage{pdflscape} % quere Seiten ermöglichen


%%%%%%%%%%%%%%% Testing %%%%%%%%%%%%%%%%%
\usepackage[format=plain,justification=centerlast]{caption} % centered captions (graphs)

\usepackage{colortbl} % Wichtig für farbige Linien

\sloppy
\interfootnotelinepenalty=10000
\DeclareUnicodeCharacter{0308}{\"}
\usepackage{romannum}
\usepackage{siunitx} % Degree celsius \SI{29}{\degreeCelsius}

\usepackage{multirow}
%\usepackage{microtype} % against hbox under/overfull

